\section*{Série 2 -- Exercice 3}

Soit $f$ la fonction définie sur $\mathbb{R}$ par :
$$f(x) = (x-1)e^x + 1$$

\begin{enumerate}
\item Calculer $\lim_{x \to -\infty} f(x)$ et $\lim_{x \to +\infty} f(x)$.

\item
\begin{enumerate}
\item Montrer que pour tout $x \in \mathbb{R}$ : $f'(x) = x e^x$.
\item Dresser le tableau de variations de $f$.
\end{enumerate}

\item
\begin{enumerate}
\item Montrer que l'équation $f(x) = 0$ admet une unique solution $\alpha$ dans $\mathbb{R}$
et que $\alpha \in [-1, 0]$.
\item Donner un encadrement de $\alpha$ d'amplitude $0{,}25$.
\end{enumerate}

\item Soit $F$ la fonction définie sur $\mathbb{R}$ par $F(x) = (x-2)e^x + x$.
Montrer que $F$ est une primitive de $f$ sur $\mathbb{R}$.

\item Calculer l'aire $\mathcal{A}$ du domaine délimité par la courbe de $f$, l'axe des abscisses
et les droites d'équation $x = 0$ et $x = 1$.
\end{enumerate}
